\section{Related Work}

Our work sits at the intersection of three lines of research: LLM self-repair with
feedback oracles, specification grounding for code generation, and statistical design
for paired LLM repair evaluation. We organize the discussion to show where each area
falls short of the controlled comparison our design provides.

\subsection{LLM Self-Repair and Feedback Oracles}

The canonical result in LLM self-repair is Self-Refine \citep{Madaan2023}, which
shows $+$20\% absolute improvement on multiple tasks with iterative self-generated feedback.
Self-Refine's iterative self-feedback loop --- where the model generates its own critique
and refines its output across multiple rounds --- motivates our Condition B design: a
fixed reprompt that re-exposes the model to the task without providing external error
context. Condition B simplifies Self-Refine's iterative self-feedback to a single fixed
reprompt string, isolating re-exposure from self-generated feedback. The result is
compelling but does not isolate what type of feedback is necessary --- it conflates
re-exposure with informative feedback.

FeedbackEval \citep{Dai2025} is the closest systematic study of feedback type
effects on code repair. On their benchmark, mixed feedback achieves 63.6\% fix rate
vs. 53.1\% for minimal feedback. Critically, removing docstrings from feedback prompts
causes severe performance degradation --- a direct empirical motivation for our
docstring inclusion. However, FeedbackEval uses a purpose-built benchmark rather than
EvalPlus, and does not isolate the value of specification context over blind
re-exposure through a McNemar paired design.

\citet{Arimbur2026} evaluates iterative self-repair on HumanEval, finding $+$4.9 to $+$17.1pp
improvement over several rounds, with assertion errors proving hardest ($\sim$45\% failure
rate). This confirms that EvalPlus-style failures are difficult for self-repair, but
does not test structured specification context.

The DUALFIX pipeline \citep{Akli2026} demonstrates that execution-feedback repair
leaves a residual class of failures addressable only by specification-level
intervention --- evolved repair rules fix 12--17\% of cases that execution repair cannot.
This is direct support for the claim that semantic-level context is necessary, not
just error traces.

\subsection{Specification Grounding for Code Generation}

\citet{Haeri2026} provide the strongest direct support for our hypothesis.
They show specification grounding improves correct code generation by $+$38 percentage
points, with docstring context identified as the primary driver. The result is on
code generation, not code repair, but the mechanism --- docstring re-anchoring preventing
goal misspecification --- transfers directly to the repair setting.

ContrastRepair \citep{Kong2024} tests contrastive I/O pairs for automated program
repair on Defects4J, finding 143/337 bugs fixed vs. 124 with single failure messages.
The contrastive pair design --- expected output alongside actual output --- is analogous
to our I/O pair component (Condition C element 2). The key difference is that
ContrastRepair targets Java bugs in a classical APR setting, not LLM code generation
on EvalPlus, and does not include docstring re-anchoring as a third component.

VRpilot \citep{Kulsum2024} applies CoT reasoning with patch validation to
vulnerability repair, finding $+$14\% correct patches over naive repair. Structured
semantic context consistently outperforms minimal feedback across repair domains.

SGCR \citep{Wang2025} validates specification-grounded code review at industry
scale, with 42\% developer adoption rate and 90.9\% improvement in review utility.

\subsection{Statistical Design for LLM Repair Evaluation}

\citet{Iscan2026} is the methodological foundation for our McNemar design. The paper
validates placebo-controlled decomposition of self-repair feedback, showing that
content (code$+$facts, $+$18pp over bare code, p$=$0.00042) rather than re-exposure drives
repair improvement. Iscan explicitly validates the one-tailed McNemar design for
paired LLM repair comparisons --- the same design we pre-register for our mechanism
experiments.

The EvalPlus benchmark \citep{Liu2023} is the evaluation standard we build on.
Its 80$\times$ test augmentation over base HumanEval makes it a substantially stricter
evaluation environment, which is precisely why semantic failures that survive round-0
are the most challenging and interesting target for repair research.

\subsection{Our Position}

Relative to this literature, our work contributes (1) an explicit verification of the
EvalPlus failure set as a reproducible data source for repair experiments,
(2) the first 3-condition McNemar design isolating specification-aligned repair from
blind reprompting on EvalPlus semantic failures, and (3) an explicit pre-registration
of the comparative predictions before executing the mechanism experiments.
